\begin{abstract}
针对微网在负荷、光伏与电价变化下的购电调控问题，建立历史预测、两日计划、日内调整与逐槽执行相结合的优化模型。各层明确区分已知预报和未来实际值，以跨日连续储能协调当前供需与次日需求。

问题一建立包含能量平衡、充放电效率、容量及功率界限和日循环条件的线性规划。典型日最优购电量为59482.7 kWh，费用35126.95元，较无储能基线节省26.9\%。

问题二以历史实现费用标定三源凸组合预测，采用半衰期5天的指数加权更新。开发期选定负荷抬升系数1.02与光伏折减25 kW，每天优化两日、执行首日，实际日末电量自由并跨日传递。2月至12月共334天总费用为1395.7万元，其中紧急购电费79.6万元。固定计划的联合残差重采样显示，总成本均值1464万元，P95为1542万元。

问题三将官方光伏预报与历史预测按0.7和0.3组合，在6、12、18时调整未执行计划，并在前两次调整引入40个联合残差情景。主方案费用1349.2万元，较官方预报无调整基线下降7.87\%。消融中无对冲方案费用1347.1万元，表明当前数据支持日内调整和组合预测的价值，但尚不支持对冲的实现费用优势。

问题四将价格替换为给定波动电价，两层策略费用分别为1459.8和1412.6万元。仅用历史价格预测的扩展费用为1465.9和1420.5万元，对应额外成本0.42\%和0.56\%。

情景数、种子与历史池扰动下结果总体稳定，约束残差达到浮点误差量级，跨日SOC连续且场景池无前视。模型同时说明初始状态、单年样本与分层规划近似的适用边界，为微网购电和储能运行提供可复现的决策依据。
\keywords{微网调度；两日滚动；因果预测；储能优化；日内调整；联合残差情景}
\end{abstract}
